\documentclass[11pt]{amsart}

\usepackage[T1]{fontenc}
\usepackage{lmodern}
\usepackage{microtype}
\usepackage{amsmath,amssymb,amsthm}
\usepackage[hidelinks]{hyperref}

\hypersetup{
  pdftitle={A Steiner Quadruple System of Order 38 with Minimum-Colorable Derived Designs}
}

\newtheorem{theorem}{Theorem}
\newtheorem{proposition}[theorem]{Proposition}
\newtheorem{lemma}[theorem]{Lemma}
\newtheorem{corollary}[theorem]{Corollary}

\newcommand{\Z}{\mathbb Z}
\newcommand{\STS}{\mathrm{STS}}
\newcommand{\SQS}{\mathrm{SQS}}
\newcommand{\mcDSQS}{\mathrm{mcDSQS}}
\newcommand{\Orb}{\operatorname{Orb}}

\title[An \texorpdfstring{$\mcDSQS(38)$}{mcDSQS(38)}]
{A Steiner Quadruple System of Order 38 with Minimum-Colorable Derived Designs}

\date{}

\subjclass[2020]{Primary 05B05; Secondary 05B07, 94B65}
\keywords{Steiner quadruple system, Steiner triple system, block coloring,
partial parallel class, derived design}

\begin{document}

\begin{abstract}
We construct a Steiner quadruple system of order $38$ whose derived
Steiner triple system at every point has block-chromatic index $19$.
Since every $\STS(37)$ requires at least $19$ partial parallel classes,
this is an $\mcDSQS(38)$, resolving the smallest unknown
$\mcDSQS(6n+2)$ case posed by Tan and Zhou. It also establishes the
$n=38$ case of Conjecture~2 of Shi, Xia, and Krotov and yields
$q^{\prime}_{0}(3,4,38)=20$. The system is specified by $27$ orbits
under an affine group of order $111$, and the two colorings that the
proof needs are printed in full.
\end{abstract}

\maketitle

\section{Result}

A block coloring of an $\STS(v)$ is a partition of its blocks into
partial parallel classes. For an $\STS(v)$ $\mathcal D$, let
$\chi'(\mathcal D)$ be the least number of classes in such a coloring,
and set
\[
\chi'(v)=\min\{\chi'(\mathcal D):\mathcal D\text{ is an }\STS(v)\}.
\]
Following Tan and Zhou~\cite{TanZhou}, an $\SQS(v)$ is an
$\mcDSQS(v)$ if every derived $\STS(v-1)$ has chromatic index
$\chi'(v-1)$.

\begin{theorem}\label{thm:main}
There exists an $\mcDSQS(38)$. More precisely, there is an $\SQS(38)$
for which every derived $\STS(37)$ has block-chromatic index $19$.
\end{theorem}

Tan and Zhou asked for an $\mcDSQS(38)$ and identified it as the
smallest unknown $\mcDSQS(6n+2)$~\cite[Sec.~4, item~(1)]{TanZhou}.
Theorem~\ref{thm:main} answers that part of their question; it does not
address the $\mathrm{RDSQS}(52)$ requested in the same item. It also
establishes the $n=38$ instance of~\cite[Conj.~2]{ShiXiaKrotov}.
Rotational $\SQS(38)$s were already known: Feng, Chang, and Ji
constructed rotational $\SQS(p+1)$s with a nontrivial multiplier for
primes $p\equiv13\pmod{24}$~\cite{FengChangJi}. The value
$\chi'(37)=19$ is known as well: Tan and Zhou record
$\chi'(v)=(v+1)/2$ for $v\equiv1\pmod{6}$ with
$v\notin\{7,13\}$~\cite{TanZhou}. The contribution here
is the simultaneous minimum colorability of all $38$ derived systems.

\section{Construction}

Let
\[
X=\Z_{37}\cup\{\infty\}.
\]
All arithmetic on finite points is modulo $37$. Since $10$ has order
$3$ in $\Z_{37}^{\times}$, the maps
\[
g_{a,b}(x)=10^a x+b,
\qquad a\in\{0,1,2\},\quad b\in\Z_{37},
\]
fixing $\infty$, form a group $G$ of order $111$. For $Q\subseteq X$,
write $\Orb_G(Q)=\{g(Q):g\in G\}$.

Consider the representatives
\[
\mathcal A=\left\{
\begin{array}{lll}
\{\infty,0,1,11\},&
\{\infty,0,2,22\},&
\{\infty,0,4,7\},\\
\{\infty,0,8,14\},&
\{\infty,0,16,28\},&
\{\infty,0,19,32\}
\end{array}
\right\},
\]
\[
\mathcal S=\left\{
\begin{array}{lll}
\{0,1,27,34\},&
\{0,2,17,31\},&
\{0,3,7,28\},\\
\{0,2,19,24\},&
\{0,5,23,29\},&
\{0,1,12,28\}
\end{array}
\right\},
\]
and
\[
\mathcal L=\left\{
\begin{array}{lll}
\{0,1,13,14\},&
\{0,1,8,21\},&
\{0,1,9,20\},\\
\{0,1,10,31\},&
\{0,1,16,23\},&
\{0,1,2,7\},\\
\{0,1,5,33\},&
\{0,1,4,18\},&
\{0,1,17,30\},\\
\{0,1,3,15\},&
\{0,1,19,35\},&
\{0,2,9,15\},\\
\{0,2,8,33\},&
\{0,2,18,23\},&
\{0,2,25,32\}
\end{array}
\right\}.
\]
Define
\[
\mathcal B=
\bigcup_{Q\in\mathcal A\cup\mathcal S\cup\mathcal L}\Orb_G(Q).
\]

\begin{proposition}\label{prop:sqs}
The pair $(X,\mathcal B)$ is an $\SQS(38)$.
\end{proposition}

\begin{proof}
If $g_{a,b}$ fixes a block with $k$ finite points of sum $\sigma$,
then summing those points gives $kb=(1-10^{a})\sigma$ in $\Z_{37}$, so
$b$ is determined by $a$; as $|G|=111$, each stabilizer therefore has
order $1$ or $3$. Testing the two candidates for each representative shows that
each representative in $\mathcal A\cup\mathcal S$ has stabilizer of
order $3$, generated by $g_{1,b}$ with
\[
b=1,2,4,8,16,32\ \text{ on }\mathcal A,
\qquad
b=27,17,7,19,29,28\ \text{ on }\mathcal S
\]
in the order listed, and that each representative in $\mathcal L$ has
trivial stabilizer. A direct computation confirms that the $27$ orbits
are pairwise disjoint. Thus
\[
|\mathcal B|=12\cdot37+15\cdot111=2109.
\]
Enumerating the triples contained in these blocks gives
\[
\binom{38}{3}=8436
\]
distinct triples, each with multiplicity one. Hence every triple of
$X$ lies in a unique block of $\mathcal B$.
\end{proof}

For $x\in X$, let
\[
\mathcal D_x=
\bigl(X\setminus\{x\},
\{B\setminus\{x\}:x\in B\in\mathcal B\}\bigr)
\]
be the design derived at $x$. Proposition~\ref{prop:sqs} implies that
$\mathcal D_x$ is an $\STS(37)$ for every $x\in X$.

\section{Minimum colorability}

\begin{lemma}\label{lem:lower}
Every $\STS(37)$ has block-chromatic index at least $19$.
\end{lemma}

\begin{proof}
An $\STS(37)$ has $37\cdot36/6=222$ blocks, whereas a partial parallel
class contains at most $\lfloor37/3\rfloor=12$ blocks. Therefore every
block coloring has at least
\[
\left\lceil\frac{222}{12}\right\rceil=19
\]
classes.
\end{proof}

The action of $G$ has the two point orbits $\{\infty\}$ and
$\Z_{37}$. The multiplier $\mu=g_{1,0}$ fixes $\infty$ and $0$, so it
is an automorphism of both $\mathcal D_{\infty}$ and $\mathcal D_0$.
Tables~\ref{tab:inf} and~\ref{tab:zero} list, for each of these two
designs, a set $C_0$ of ten blocks and six sets $G_1,\dots,G_6$ of
twelve blocks each; the classes are $C_0$ together with the eighteen
sets
\[
\mu^{j}(G_i)\setminus C_0
\qquad(1\leq i\leq6,\ 0\leq j\leq2).
\]
This partitions each of $\mathcal D_{\infty}$ and $\mathcal D_0$ into
$19$ partial parallel classes, in both cases with class-size profile
\[
12^{14},11^4,10;
\]
that is, fourteen classes of size $12$, four of size $11$, and one of
size $10$. Every listed triple is a block of the relevant derived
design, the triples within each class are pairwise disjoint, and the
classes partition all $222$ blocks.

\begin{table}[htb]
\caption{A $19$-coloring of $\mathcal D_{\infty}$, whose point set is $\Z_{37}$.}\label{tab:inf}
\footnotesize
\begin{tabular}{@{}r@{\;\;}l@{}}
$C_0$ & $\{1,10,26\}$\;$\{2,15,20\}$\;$\{3,11,17\}$\;$\{5,13,19\}$\;$\{6,18,27\}$\;$\{7,33,34\}$ \\
 & $\{8,16,22\}$\;$\{9,35,36\}$\;$\{14,29,31\}$\;$\{21,25,28\}$ \\
$G_1$ & $\{0,30,34\}$\;$\{2,8,31\}$\;$\{3,11,17\}$\;$\{4,5,15\}$\;$\{6,19,24\}$\;$\{7,9,29\}$ \\
 & $\{10,12,32\}$\;$\{13,21,27\}$\;$\{14,22,28\}$\;$\{16,20,23\}$\;$\{18,33,35\}$\;$\{25,26,36\}$ \\
$G_2$ & $\{0,5,24\}$\;$\{1,17,29\}$\;$\{3,8,27\}$\;$\{4,16,25\}$\;$\{6,10,13\}$\;$\{7,12,31\}$ \\
 & $\{9,35,36\}$\;$\{11,23,32\}$\;$\{14,18,21\}$\;$\{15,19,22\}$\;$\{20,28,34\}$\;$\{26,30,33\}$ \\
$G_3$ & $\{0,23,31\}$\;$\{1,6,25\}$\;$\{2,28,29\}$\;$\{3,26,34\}$\;$\{4,19,21\}$\;$\{7,13,36\}$ \\
 & $\{8,16,22\}$\;$\{9,10,20\}$\;$\{11,15,18\}$\;$\{12,24,33\}$\;$\{14,27,32\}$\;$\{17,30,35\}$ \\
$G_4$ & $\{0,1,11\}$\;$\{2,18,30\}$\;$\{3,16,21\}$\;$\{4,6,26\}$\;$\{5,31,32\}$\;$\{8,20,29\}$ \\
 & $\{9,22,27\}$\;$\{10,14,17\}$\;$\{12,13,23\}$\;$\{15,28,33\}$\;$\{19,34,36\}$\;$\{24,25,35\}$ \\
$G_5$ & $\{0,16,28\}$\;$\{1,24,32\}$\;$\{2,7,26\}$\;$\{3,29,30\}$\;$\{4,12,18\}$\;$\{5,20,22\}$ \\
 & $\{6,21,23\}$\;$\{8,17,33\}$\;$\{9,11,31\}$\;$\{10,19,35\}$\;$\{13,25,34\}$\;$\{15,27,36\}$ \\
$G_6$ & $\{0,2,22\}$\;$\{1,9,15\}$\;$\{3,18,20\}$\;$\{4,8,11\}$\;$\{5,14,30\}$\;$\{6,7,17\}$ \\
 & $\{10,25,27\}$\;$\{12,16,19\}$\;$\{13,26,31\}$\;$\{23,24,34\}$\;$\{28,32,35\}$\;$\{29,33,36\}$ \\
\end{tabular}
\end{table}

\begin{table}[htb]
\caption{A $19$-coloring of $\mathcal D_{0}$, whose point set is $(\Z_{37}\setminus\{0\})\cup\{\infty\}$.}\label{tab:zero}
\footnotesize
\begin{tabular}{@{}r@{\;\;}l@{}}
$C_0$ & $\{1,5,33\}$\;$\{3,4,30\}$\;$\{6,8,23\}$\;$\{7,19,26\}$\;$\{9,12,16\}$\;$\{10,13,34\}$ \\
 & $\{11,27,36\}$\;$\{14,21,29\}$\;$\{17,22,35\}$\;$\{18,24,32\}$ \\
$G_1$ & $\{1,16,23\}$\;$\{2,14,36\}$\;$\{4,6,12\}$\;$\{5,18,20\}$\;$\{7,19,26\}$\;$\{8,13,31\}$ \\
 & $\{9,22,28\}$\;$\{10,24,35\}$\;$\{11,21,25\}$\;$\{15,27,33\}$\;$\{17,29,32\}$\;$\{30,34,\infty\}$ \\
$G_2$ & $\{1,2,7\}$\;$\{3,5,14\}$\;$\{4,20,26\}$\;$\{8,18,27\}$\;$\{9,10,21\}$\;$\{11,24,33\}$ \\
 & $\{12,34,35\}$\;$\{13,19,28\}$\;$\{15,17,\infty\}$\;$\{16,29,36\}$\;$\{22,23,25\}$\;$\{30,31,32\}$ \\
$G_3$ & $\{1,10,31\}$\;$\{2,8,33\}$\;$\{3,6,22\}$\;$\{4,16,27\}$\;$\{5,35,36\}$\;$\{7,18,25\}$ \\
 & $\{11,13,23\}$\;$\{12,15,20\}$\;$\{14,21,29\}$\;$\{17,28,34\}$\;$\{19,32,\infty\}$\;$\{24,26,30\}$ \\
$G_4$ & $\{1,24,25\}$\;$\{2,6,13\}$\;$\{3,19,34\}$\;$\{4,21,23\}$\;$\{5,10,17\}$\;$\{7,12,14\}$ \\
 & $\{8,20,32\}$\;$\{9,18,31\}$\;$\{15,29,35\}$\;$\{16,28,\infty\}$\;$\{22,27,30\}$\;$\{26,33,36\}$ \\
$G_5$ & $\{1,9,20\}$\;$\{2,27,35\}$\;$\{3,13,16\}$\;$\{4,10,22\}$\;$\{5,12,32\}$\;$\{6,29,\infty\}$ \\
 & $\{7,8,24\}$\;$\{11,18,30\}$\;$\{14,33,34\}$\;$\{15,19,25\}$\;$\{21,31,36\}$\;$\{23,26,28\}$ \\
$G_6$ & $\{1,11,\infty\}$\;$\{2,3,21\}$\;$\{4,25,34\}$\;$\{5,26,31\}$\;$\{6,7,16\}$\;$\{8,22,29\}$ \\
 & $\{9,32,35\}$\;$\{10,23,27\}$\;$\{12,13,36\}$\;$\{14,15,30\}$\;$\{17,19,33\}$\;$\{20,24,28\}$ \\
\end{tabular}
\end{table}

For $p\in\Z_{37}$, translation by $p$ belongs to $G$ and sends
$\mathcal D_0$ to $\mathcal D_p$. It therefore transports the coloring
at $0$ to every finite point.

\begin{proof}[Proof of Theorem~\ref{thm:main}]
Proposition~\ref{prop:sqs} gives an $\SQS(38)$. The two tables and
translation give a $19$-coloring of every derived design, while
Lemma~\ref{lem:lower} rules out fewer than $19$ colors. Hence
\[
\chi'(\mathcal D_x)=19\qquad(x\in X).
\]
The construction also shows $\chi'(37)\leq19$, and
Lemma~\ref{lem:lower} gives the reverse inequality. Thus
$\chi'(37)=19$, so every $\mathcal D_x$ is minimum colorable.
\end{proof}

\begin{corollary}\label{cor:q0}
Conjecture~2 of Shi, Xia, and Krotov holds for $n=38$, and
\[
q^{\prime}_{0}(3,4,38)=20.
\]
\end{corollary}

\begin{proof}
For an $\STS$, two vertices of the minimum-distance graph of
Shi, Xia, and Krotov are adjacent exactly when the corresponding
triples meet. Since two distinct Steiner triples meet in at most one
point, the independent sets of this graph are precisely partial
parallel classes. Their chromatic number is therefore the
block-chromatic index used here. By~\cite[Thm.~5]{ShiXiaKrotov},
\[
q^{\prime}_{0}(3,4,38)
=1+\min_{\mathcal S\text{ an }\SQS(38)}
\max_{\mathcal D\in\mathcal S'}\chi'(\mathcal D),
\]
where $\mathcal S'$ is the set of derived systems of $\mathcal S$.
Lemma~\ref{lem:lower} makes the quantity after $1$ at least $19$, and
Theorem~\ref{thm:main} attains $19$. The same theorem is exactly the
$n=38$ assertion of~\cite[Conj.~2]{ShiXiaKrotov}.
\end{proof}

\begin{thebibliography}{9}

\bibitem{FengChangJi}
T.~Feng, Y.~Chang, and L.~Ji,
\emph{Constructions for rotational Steiner quadruple systems},
J. Combin. Des. \textbf{17} (2009), no.~5, 353--368.
\href{https://doi.org/10.1002/jcd.20202}{doi:10.1002/jcd.20202}.

\bibitem{ShiXiaKrotov}
M.~Shi, Y.~Xia, and D.~S. Krotov,
\emph{A family of diameter perfect constant-weight codes from
Steiner systems},
J. Combin. Theory Ser. A \textbf{200} (2023), 105790.
\href{https://doi.org/10.1016/j.jcta.2023.105790}
{doi:10.1016/j.jcta.2023.105790}.

\bibitem{TanZhou}
Y.~Tan and J.~Zhou,
\emph{Steiner quadruple systems with minimum colorable derived designs:
constructions and applications},
Des. Codes Cryptogr. \textbf{93} (2025), no.~10, 4117--4140.
\href{https://doi.org/10.1007/s10623-025-01677-x}
{doi:10.1007/s10623-025-01677-x}.

\end{thebibliography}

\end{document}